\documentclass[12pt,a4paper]{article}

% Packages
\usepackage[utf8]{inputenc}
\usepackage[T1]{fontenc}
\usepackage{amsmath,amssymb,amsthm}
\usepackage{graphicx}
\usepackage{hyperref}
\usepackage{booktabs}
\usepackage{geometry}
\usepackage{xcolor}
\usepackage{caption}
\usepackage{subcaption}
\usepackage{enumitem}
\usepackage{url}
\usepackage{doi}

\geometry{margin=1in}

% Theorem environments
\newtheorem{definition}{Definition}
\newtheorem{proposition}{Proposition}
\newtheorem{remark}{Remark}

% Title information
\title{\textbf{Temporal Emergence from Dynamical Observer Coupling:\\
       A Matrix Model for Observer-Dependent Symmetry Breaking}\\[0.5cm]
       \large Version 5: Complete Scaling Analysis and Open Questions}

\author{Sutipong Chanpengpad\\
\small Independent Researcher, Chiang Rai, Thailand\\
\small \texttt{dhammawatthumpra@gmail.com}\\
\small \href{https://orcid.org/0009-0001-4069-8576}{ORCID: 0009-0001-4069-8576}}

\date{\today}

\begin{document}

\maketitle

\begin{abstract}
The Structural Time Framework (STF) proposes that time emerges from 
observer-dependent ordering of pre-temporal states. We present a 
computational demonstration using a matrix model where an explicit 
observer system $Y$ dynamically creates the temporal direction in 
the main system $X$, without any external parameters.

Our model employs a directional coupling mechanism where the observer 
generates a direction vector $v$ from its matrix traces, and dimensions 
of $X$ aligned with $v$ receive a temporal expansion term 
$-g_{XY} \hat{v}_\mu^2 \text{Tr}(X_\mu^4)$. Monte Carlo simulations with statistical analysis 
(5 seeds per configuration) show:

\begin{enumerate}[leftmargin=*]
\item Temporal emergence occurs robustly for observer dimensions
      $d \in \{2,3,4,5\}$, with mean ratios clustering in $4.6$--$4.8$ and
      no statistically distinguishable differences between dimensions.
\item Perfect alignment between the temporal direction $X_{\max}$ and 
      observer direction $v_{\max}$ in healthy emergence cases.
\item Two-regime scaling behavior: robust emergence for $N \leq 6$ with 
      fixed coupling, but parameter tuning required for $N \geq 7$.
\item A paired comparison between $d=3$ and $d=4$ found no significant
      difference ($t=1.043$, $p=0.356$); with $n=5$ seeds per condition
      this test has limited power and does not by itself establish
      equivalence or robustness.
\end{enumerate}

\textbf{Important Caveats:} This work represents a toy-level 
mathematical exploration using small matrix sizes ($N \leq 8$). 
Results demonstrate mathematical consistency, not physical reality.

\medskip
\noindent\textbf{Keywords:} dynamical observer, temporal emergence, 
matrix models, symmetry breaking, computational demonstration
\end{abstract}

\tableofcontents
\newpage

%=============================================================================
\section{Introduction}\label{sec:intro}
%=============================================================================

The nature of time has been debated by philosophers and physicists for 
centuries. The Structural Time Framework (STF)~\cite{chanpengpad_stf} 
proposes that time is not fundamental but emerges from observer-dependent 
ordering of pre-temporal states.

In our previous work~\cite{chanpengpad_sgoed_v2}, we presented a toy 
model where temporal emergence was driven by an external adapter parameter 
$\kappa_2$. While successful, this approach had a fundamental limitation: 
the observer was represented as a parameter, not as a dynamical system.

This paper addresses that limitation by introducing an \textbf{explicit 
dynamical observer}---a matrix system $Y$ that dynamically creates the 
temporal direction in the main system $X$ through coupling, without any 
external parameters.

\textbf{Important Disclaimer:} We emphasize that this work is a 
\textbf{toy-level computational demonstration}, not empirical physics. 
Our simulations use small matrix sizes ($N \leq 8$), finite Monte Carlo 
samples, and idealized actions. The results demonstrate mathematical 
consistency, not physical reality.

%=============================================================================
\section{Dynamical Observer Model}\label{sec:model}
%=============================================================================

We consider two coupled matrix systems:

\textbf{Main System $X$:} $D$ Hermitian matrices $X^\mu$ of size $N \times N$.

\textbf{Observer System $Y$:} $d$ Hermitian matrices $Y^a$ of size $N \times N$.

The total action is:
\begin{equation}
S_{\text{total}} = S_X + S_Y + S_{\text{coupling}}
\end{equation}

where $S_X$ and $S_Y$ are IKKT-like actions with stability gates, and the 
directional coupling is:
\begin{equation}
S_{\text{coupling}} = -g_{XY} \sum_\mu \hat{v}_\mu^2 \text{Tr}(X_\mu^4)
\end{equation}

The observer generates direction $\hat{v}$ from its traces. This coupling 
temporalizes dimensions aligned with $\hat{v}$.

\subsection{Theoretical Motivation for the Coupling}

The coupling term can be understood from three perspectives:

\begin{enumerate}[leftmargin=*]
\item \textbf{Effective Field Theory:} Leading-order term in $1/N$ expansion.
\item \textbf{Symmetry Breaking:} Explicitly breaks SO(D) $\rightarrow$ SO(D-1).
\item \textbf{Measurement Analogy:} Observer selects direction, system responds.
\end{enumerate}

\subsection{Why the Observer Does Not Temporalize Itself}

The distinction between observer and system is relational, not fundamental. 
The asymmetry in coupling reflects the measurement paradigm: the apparatus 
measures the system, not vice versa. This avoids the self-observation paradox 
analogous to Wigner's friend.

%=============================================================================
\section{Computational Methodology}\label{sec:method}
%=============================================================================

We use Metropolis-Hastings sampling with:
\begin{itemize}
\item Matrix size: $N \in \{4, 5, 6, 7, 8\}$
\item System dimensions: $D = 6$
\item Observer dimensions: $d \in \{2, 3, 4, 5\}$
\item Coupling strength: $g_{XY} \in [0.8, 1.2]$
\item Thermalization: 20--25 sweeps
\item Measurement: 30--40 sweeps
\item Multiple seeds: $\{42, 43, 44, 45, 46\}$
\end{itemize}

Observables measured, with $\text{extent}(X_\mu) \equiv \text{Tr}(X_\mu^2)/N$:
\begin{align}
\text{Ratio} &= \frac{\text{extent}(X_{\max})}{\text{mean}(\text{extent}(X_{\neq \max}))} \\
\text{Alignment} &= \mathbb{I}(\arg\max_\mu \text{extent}(X_\mu) == \arg\max_\mu \hat{v}_\mu^2) \\
H &= -\sum_\mu p_\mu \ln p_\mu, \quad p_\mu = \frac{\text{extent}(X_\mu)}{\sum_\nu \text{extent}(X_\nu)}
\end{align}

A configuration is classified as \textbf{Healthy} if the mean Ratio over
the five seeds satisfies $1.5 < \overline{R} < 10.0$; the upper bound
$R=10$ coincides with the stability-gate cutoff (\texttt{max\_extent}) used
in the action itself, beyond which a quadratic penalty term dominates and
suppresses further extent growth. Configurations with $\overline{R} \leq 1.5$
are labeled Weak; none of the configurations reported below exceed the
upper bound.

%=============================================================================
\section{Results}\label{sec:results}
%=============================================================================

\subsection{Statistical Robustness}

Table~\ref{tab:stats} summarizes results for optimal parameters.

\begin{table}[ht]
\centering
\caption{Statistical analysis ($N=6$, $g_{XY}=0.8$)}
\label{tab:stats}
\begin{tabular}{ccccc}
\toprule
$d$ & Mean Ratio & Std Dev & CV (\%) & Alignment \\
\midrule
3 & 4.82 & $\pm$0.24 & 5.0\% & 100\% (5/5) \\
4 & 4.60 & $\pm$0.28 & 6.0\% & 100\% (5/5) \\
\bottomrule
\end{tabular}
\end{table}

Paired t-test: $t = 1.043$, $p = 0.356$. This fails to reject the null
of no difference between $d=3$ and $d=4$; note that with $n=5$ seeds per
condition, absence of a significant difference is not evidence of
equivalence, and this test alone does not establish robustness.

%=============================================================================
\subsection{Observer-Dimension Phase Diagram}\label{sec:phase_diagram}
%=============================================================================

To investigate the role of observer complexity, we varied the observer
dimension $d \in \{2,3,4,5\}$ while keeping $N=6$, $D=6$, and
$g_{XY}=0.8$ fixed. Each configuration was evaluated over five independent
random seeds.

\begin{table}[ht]
\centering
\caption{Phase diagram by observer dimension ($N=6$, $g_{XY}=0.8$)}
\label{tab:phase_diagram}
\begin{tabular}{cccccc}
\toprule
$d$ & Mean Ratio & Std Dev & CV (\%) & Alignment & Status \\
\midrule
2 & 4.81 & $\pm$0.46 & 9.5\% & 100\% (5/5) & Healthy \\
3 & 4.82 & $\pm$0.24 & 5.0\% & 100\% (5/5) & Healthy \\
4 & 4.60 & $\pm$0.28 & 6.0\% & 100\% (5/5) & Healthy \\
5 & 4.62 & $\pm$0.16 & 3.4\% & 100\% (5/5) & Healthy \\
\bottomrule
\end{tabular}
\end{table}

Several features are noteworthy.

\begin{enumerate}[leftmargin=*]
\item \textbf{No strict threshold at $d=3$:}
Contrary to our initial expectation, even $d=2$ produces healthy temporal
emergence with perfect alignment. This suggests that the minimum observer
complexity required for temporalization is lower than anticipated.

\item \textbf{No significant differences across $d$:}
Mean ratios cluster in a narrow range ($4.6$--$4.8$) across all four tested
observer dimensions, and alignment is perfect (100\%, 5/5 seeds) in every
case. CV varies between $3.4\%$ and $9.5\%$ with no clear monotonic trend,
and the paired comparison between $d=3$ and $d=4$ (Table~\ref{tab:stats})
found no significant difference. We do not find evidence for an optimal or
``sweet spot'' observer dimension in this range; within the resolution of
five seeds per configuration, the mechanism appears robust to the choice of
$d \in \{2,3,4,5\}$. (An earlier draft of this table, generated before the
simulation pipeline was consolidated into \texttt{code/sgoed\_core.py}, had
reported a spurious $d=4$ ``sweet spot'' and reduced $d=5$ alignment; those
numbers could not be reproduced from the current codebase and have been
replaced with freshly re-run, seed-logged values.)
\end{enumerate}

%=============================================================================
\subsection{Finite-Size Scaling and Parameter Sensitivity}\label{sec:fss}
%=============================================================================

We next tested whether the observed temporal emergence is a finite-size
artifact by varying the matrix size $N$. We first fixed $d=3$ and
$g_{XY}=0.8$ and examined $N \in \{4,5,6,7,8\}$.

\begin{table}[ht]
\centering
\caption{Finite-size scaling with fixed $g_{XY}=0.8$ ($d=3$)}
\label{tab:fss_fixed}
\begin{tabular}{cccccc}
\toprule
$N$ & Mean Ratio & Std Dev & CV (\%) & Alignment & Status \\
\midrule
4 & 4.87 & $\pm$0.63 & 12.9\% & 100\% (5/5) & Healthy \\
5 & 4.90 & $\pm$0.44 & 9.0\%  & 100\% (5/5) & Healthy \\
6 & 4.82 & $\pm$0.24 & 5.0\%  & 100\% (5/5) & Healthy \\
7 & 1.83 & $\pm$0.42 & 23.0\% & 80\% (4/5)  & Weak \\
8 & 1.70 & $\pm$0.50 & 29.4\% & 100\% (5/5) & Weak \\
\bottomrule
\end{tabular}
\end{table}

The data reveal two distinct regimes.

\begin{enumerate}[leftmargin=*]
\item \textbf{Small-$N$ regime ($N \leq 6$):}
For $N=4,5,6$, the ratio remains stable near $R \approx 4.8$--$4.9$,
with perfect alignment. This indicates that the mechanism is robust in
the small-$N$ regime.

\item \textbf{Crossover regime ($N \approx 7$):}
At $N=7$, the ratio drops to $R \approx 1.83$ and the alignment rate falls
to 80\%. This suggests the onset of a crossover where the fixed coupling
$g_{XY}=0.8$ is no longer sufficient to produce strong temporalization.

\item \textbf{Large-$N$ sensitivity ($N \geq 7$):}
At $N=8$, emergence remains weak for fixed $g_{XY}=0.8$, despite perfect
alignment. This indicates that the observer direction is still correctly
selected, but the coupling strength is insufficient to amplify it into a
strong temporal hierarchy.
\end{enumerate}

%=============================================================================
\subsection{Parameter Tuning at $N=8$}\label{sec:n8_tuning}
%=============================================================================

To determine whether the weak emergence at $N=8$ represents a genuine
breakdown or merely parameter sensitivity, we repeated the simulations at
$N=8$ while varying $g_{XY}$.

\begin{table}[ht]
\centering
\caption{Parameter tuning at $N=8$ ($d=3$)}
\label{tab:n8_tuning}
\begin{tabular}{cccccc}
\toprule
$g_{XY}$ & Mean Ratio & Std Dev & CV (\%) & Alignment & Status \\
\midrule
0.80 & 1.70 & $\pm$0.50 & 29.4\% & 100\% (5/5) & Weak \\
1.05 & 5.72 & $\pm$3.37 & 59.0\% & 100\% (5/5) & High variance \\
1.10 & 7.41 & $\pm$2.33 & 31.5\% & 100\% (5/5) & Healthy / preferred \\
1.15 & 8.45 & $\pm$2.24 & 26.5\% & 100\% (5/5) & Near boundary \\
\bottomrule
\end{tabular}
\end{table}

These results show that the apparent failure at $N=8$ is not a genuine
absence of the mechanism. Instead, $N=8$ requires stronger coupling to
restore temporal emergence.

The value $g_{XY}=1.10$ provides a useful compromise: it produces a healthy
mean ratio $R \approx 7.41$, perfect alignment, and remains safely below the
upper boundary $R=10$ on average. We therefore regard
\[
g_{XY}^{\mathrm{opt}}(N=8) \approx 1.10
\]
as the preferred working value in this parameter range.

We note that CV is non-monotonic in $g_{XY}$ over the scanned range,
peaking at $g_{XY}=1.05$ (CV $=59.0\%$) before falling again at
$g_{XY}=1.10, 1.15$. The cause of this spike is unclear from the present
data: with $n=5$ seeds at a single point, we cannot distinguish an
ordinary sampling fluctuation from a genuine feature of the model near
$g_{XY}=1.05$. We flag this as an open question (Section~\ref{sec:open_questions})
rather than propose a mechanism for it.

%=============================================================================
\subsection{Two-Regime Scaling Picture}\label{sec:two_regime_scaling}
%=============================================================================

The combined finite-size and tuning data suggest a two-regime scaling
structure.

\begin{table}[ht]
\centering
\caption{Estimated optimal coupling by system size}
\label{tab:gopt_scaling}
\begin{tabular}{ccccc}
\toprule
$N$ & $g_{XY}^{\mathrm{opt}}$ & Mean Ratio & Alignment & Regime \\
\midrule
4 & 0.80 & 4.87 & 100\% & Small-$N$ \\
5 & 0.80 & 4.90 & 100\% & Small-$N$ \\
6 & 0.80 & 4.82 & 100\% & Small-$N$ \\
7 & $>0.80$ & 1.83 at $g=0.8$ & 80\% & Crossover \\
8 & $\approx 1.10$ & 7.41 & 100\% & Tuned large-$N$ \\
\bottomrule
\end{tabular}
\end{table}

For $N \leq 6$, the same coupling $g_{XY}=0.8$ produces stable temporal
emergence. For $N \geq 7$, however, the system enters a regime where the
coupling must be increased to maintain a strong temporal hierarchy.

We caution that only two anchor values of $g_{XY}^{\mathrm{opt}}$ are
available ($\approx 0.8$ for $N\le 6$, and $\approx 1.10$ at $N=8$, itself
selected by eye from four scanned values rather than fit), so any
``power-law exponent'' computed from them is a secant slope between two
points, not an independent statistical estimate, and the two numbers below
are not separate pieces of evidence but different pairings of the same two
anchors. If one nonetheless assumes a power-law form
\[
g_{XY}^{\mathrm{opt}}(N) \propto N^\alpha,
\]
then using the range $N=6 \rightarrow 8$ gives
\[
\alpha
= \frac{\ln(1.10/0.80)}{\ln(8/6)}
\approx
\frac{0.318}{0.288}
\approx 1.10,
\]
while using the wider range $N=4 \rightarrow 8$ (which additionally assumes
no change in $g_{XY}^{\mathrm{opt}}$ between $N=4$ and $N=6$) gives
\[
\alpha
= \frac{\ln(1.10/0.80)}{\ln(8/4)}
\approx
\frac{0.318}{0.693}
\approx 0.46.
\]
The factor-of-two spread between these two secant slopes is expected from
their construction and should not be read as evidence against a power law;
conversely, agreement between them would not have been evidence for one
either, since both are computed from the same two data points. With only
two anchors, the data cannot distinguish a power law from other
monotonically increasing forms of $g_{XY}^{\mathrm{opt}}(N)$. Determining
$\alpha$ (or ruling out a power law) would require $g_{XY}^{\mathrm{opt}}$
at intermediate points such as $N=7$ and beyond $N=8$. We therefore report
these numbers only as a preliminary, non-independent sensitivity check, not
as a scaling-exponent measurement, and the data are more naturally read as
evidence for a crossover near $N \approx 7$ between a small-$N$ regime,
where no tuning is required, and a larger-$N$ regime, where $g_{XY}$ must
increase with system size.

%=============================================================================
\subsection{Physical Interpretation of the Crossover}\label{sec:crossover_interpretation}
%=============================================================================

The observed crossover can be interpreted in several complementary ways.

\begin{enumerate}[leftmargin=*]
\item \textbf{Phase-space volume:}
Each matrix contains $N^2$ real degrees of freedom before symmetry constraints.
For $N=8$, there are $d \times N^2 = 3 \times 64 = 192$ independent
matrix elements (for $d=3$). Thermal fluctuations at inverse temperature
$\beta = 1$ compete with the coupling energy $\propto g_{XY}$. For fixed
$g_{XY}$, larger $N$ means more fluctuations, requiring larger $g_{XY}$ to
maintain the symmetry-broken phase.

\item \textbf{Competition with IKKT term:}
The IKKT commutator term scales as $\sim N^2$ in the large-$N$ limit. To
overcome this and create a preferred direction, the coupling term must also
scale appropriately, suggesting $g_{XY} \propto N^\alpha$ for some $\alpha > 0$.

\item \textbf{Entropy argument:}
The number of possible configurations grows exponentially with $N$. To select
a specific temporal direction from this exponentially large space requires
coupling energy that grows with system size.
\end{enumerate}

%=============================================================================
\subsection{Connection to IKKT Matrix Model}

This scaling behavior is reminiscent of the IKKT matrix
model~\cite{kim_nishimura_2011, kim_nishimura_2019}. In the Lorentzian IKKT
model, Kim and Nishimura found that spontaneous symmetry breaking
SO(9) $\rightarrow$ SO(3) requires $N \geq 10$ for clear emergence of
(3+1)-dimensional spacetime. While the coupling in their model is fixed by
$g^2$, the large-$N$ limit requires careful extrapolation and shows
sensitivity to initial conditions.

Our model shows a similar sensitivity, but with an explicit scaling parameter
$g_{XY}$ that must be tuned with $N$. This suggests that both spontaneous and
observer-induced symmetry breaking face similar challenges in the large-$N$
limit, namely the competition between coupling energy and phase space volume.

%=============================================================================
\section{Discussion}\label{sec:discussion}
%=============================================================================

\subsection{Advantages Over External Adapter}

The dynamical observer approach offers several advantages:

\begin{enumerate}[leftmargin=*]
\item \textbf{Physical Motivation:} Observer is dynamical system, not parameter.
\item \textbf{Emergent Direction:} Temporal direction emerges from observer dynamics.
\item \textbf{Robustness:} Works across $d = 2,3,4,5$ and multiple seeds.
\item \textbf{Measurement Theory:} Observer as measurement device.
\end{enumerate}

\subsection{Relation to Philosophical Positions}

Our model is consistent with several philosophical positions within the 
restricted scope of our toy model:

\begin{itemize}[leftmargin=*]
\item \textbf{Leibniz (relational time):} Temporal direction created by observer.
\item \textbf{Barbour (timeless physics):} Pre-observer state is symmetric.
\item \textbf{Rovelli (thermal time):} Coupling creates preferred direction.
\item \textbf{Relational QM:} States relative to observer.
\end{itemize}

\textbf{Important:} ``Consistent with'' does not mean ``proves'' or 
``validates.'' These philosophical positions remain open questions.

%=============================================================================
\section{Open Questions and Future Directions}\label{sec:open_questions}
%=============================================================================

Our results raise several important questions that warrant further investigation:

\begin{enumerate}[leftmargin=*]
\item \textbf{Precise Scaling Exponent:}
Our preliminary estimate $\alpha \approx 0.46$--$1.10$ is based on limited
data points. More systematic studies at intermediate system sizes
($N=5, 7, 9, 10$) are needed to confirm the power-law scaling hypothesis
and determine whether $\alpha$ is constant or varies with $N$.

\item \textbf{Nature of the Crossover:}
Is the crossover at $N \approx 7$ a sharp transition or a smooth crossover?
Does it depend on observer dimension $d$? Would increasing $d$ at large $N$
reduce the required coupling strength?

\item \textbf{Large-$N$ Behavior:}
The scaling behavior at $N > 8$ remains unexplored. Key questions include:
Does the scaling $g_{XY}^{\text{opt}}(N) \propto N^\alpha$ persist? Is there
a critical $N_c$ where the mechanism breaks down permanently? How does the
continuum limit ($N \rightarrow \infty$) behave?

\item \textbf{Sweet Spot Determination:}
At $N=8$, the optimal coupling appears to be $g_{XY}^{\text{opt}} \approx 1.10$.
Finer sampling in this range (e.g., $g_{XY} = 1.08, 1.12$) would identify
the precise sweet spot where emergence is strongest with minimal variance.

\item \textbf{Observer Complexity Scaling:}
How does the optimal observer dimension $d^{\text{opt}}$ scale with $N$?
Is there a relation $d^{\text{opt}}(N) \propto N^\beta$? Would larger $d$
at large $N$ restore emergence without requiring stronger coupling?

\item \textbf{Universality:}
Would this mechanism work with different coupling structures or different
matrix models (e.g., BFSS, Lorentzian IKKT)? How universal is the
observer-induced temporalization phenomenon?

\item \textbf{Connection to Quantum Gravity:}
The finite-size scaling and crossover behavior may have analogs in emergent
spacetime scenarios. Can this model be connected to holographic duality or
other quantum gravity frameworks?

\item \textbf{Renormalization Group Analysis:}
The scaling behavior suggests the need for RG analysis to understand how
$g_{XY}$ flows with scale. Such an analysis would determine whether the
scaling is a physical feature of the continuum theory or an artifact of
our phenomenological coupling.

\item \textbf{Non-Monotonic Variance Near $g_{XY}=1.05$:}
At $N=8$, the coefficient of variation peaks sharply at $g_{XY}=1.05$
(CV $=59.0\%$) between lower values at $g_{XY}=0.80$ and higher values at
$g_{XY}=1.10,1.15$ (Section~\ref{sec:n8_tuning}). Is this a reproducible
feature of the model (e.g.\ proximity to a region of enhanced
fluctuations) or an artifact of $n=5$ seeds at a single point? Denser
sampling around $g_{XY}=1.05$ with more seeds would be needed to
distinguish these possibilities.
\end{enumerate}

%=============================================================================
\section{Conclusion}\label{sec:conclusion}
%=============================================================================

We have presented a toy-level computational demonstration of temporal 
emergence from a dynamical observer. Our key findings:

\begin{enumerate}[leftmargin=*]
\item An observer system $Y$ can dynamically create the temporal direction 
      in a main system $X$ through directional coupling.
\item For coupling $g_{XY} \in [0.8, 1.2]$ and observer dimensions 
      $d \in \{2, 3, 4, 5\}$, we observe healthy temporal emergence with 
      ratios $1.8$--$8.5$.
\item Perfect alignment ($X_{\max} = v_{\max}$) in healthy emergence cases 
      confirms the observer's active role.
\item Two-regime scaling: robust for $N \leq 6$, parameter tuning required 
      for $N \geq 7$.
\item A paired comparison between $d=3$ and $d=4$ found no significant
      difference ($p = 0.356$), though with $n=5$ seeds per condition
      this does not by itself establish robustness or equivalence.
\end{enumerate}

\textbf{Appropriate Interpretation:} These results demonstrate that 
dynamical observer temporalization \textit{can be represented} in a 
consistent mathematical framework. They do not prove that time emerges 
from observer coupling in nature.

\textbf{Final Statement:} This work is a toy-level demonstration of 
mathematical consistency, not a discovery of physical laws. Following 
the philosophy of our exploratory notes, we emphasize that this is a 
\textbf{design framework demonstrating mathematical consistency}, 
not empirical physics.

\bibliographystyle{plain}
\bibliography{references}

\end{document}
